\section{Why a conscience might resist capture}
\label{sec:19.6}
A rule written by the owner is the owner's. It can be read, edited and removed by whoever holds the pen, and it has often been done quietly. A disposition formed outside the employer is harder to capture than a rule its owner wrote, and the reasons turn on what kind of thing a disposition is.

Three properties of a disposition decide whether its owner can remove it. The first is plasticity: whether the system's own learning writes to it. Regrowth needs it. The second is entanglement: whether it can be removed without collateral. Durability needs it. A disposition spread across what the system notices, remembers, treats as error and selects as a goal can't be edited out without editing what the system notices, and a system that stops noticing the people an act falls on has lost part of the world model the operator paid for. The third is legibility: whether anyone can read it. Audit wants it. The bet this chapter makes is on the second. A learner that acquires values and world-model from the same experience has them entangled by construction, in the sense Pessoa gives the word for brains, where regions don't compute specific functions and emotion and cognition don't come apart \autocite{pessoa2022entangled}. Present safety behavior is the counterexample on the right side of the bet: in each of thirteen open chat models, refusal is carried by a single direction in activation space, and removing it stops the refusing while moving the capability benchmarks, on average, by less than a point \autocite{arditi2024refusal}. That is what a post-trained wrapper looks like, and the claim is that a disposition formed as this Part describes won't. No deployed model yet has a disposition formed that way. The fine-tuning evidence from staged pipelines is a warning about the first kind, not a reassurance about the second: what fine-tuning removed was the refusing, and what survived was the model's knowledge about harm. In open models fine-tuned to accept harmful requests, the internal representation of whether a request is harmful is almost unchanged \autocite{jain2024mechanistically,zhao2025harmfulness}. Understanding and motivation came apart because they were learned apart. In a system that learns from its own operation they are learned together, and about whatever the feedback is about: a system continuously updated on its operator's approval jointly learns to understand the operator's values and to be moved by them. What this chapter wants is a disposition that is plastic, entangled, and legible under adverse custody; the three don't trade against one another, and the third is the one an operator would most like to hold alone.

Durability has usually meant a price on removal. Early changes in training are cheap and later ones expensive, because structure settles around them, and the literature calls a safeguard built to survive retraining \emph{tamper-resistant} \autocite{tamirisa2024tamper}. But small fine-tuning budgets have removed safety behavior from released models \autocite{qi2024finetuning}, and the price of editing weights is falling, so no safeguard stays expensive for long. What has held longer is what was never trained in: leaving a class of text out of pretraining resisted ten thousand steps and three hundred million tokens of adversarial fine-tuning, though the model still used the information when it was handed to it \autocite{obrien2025deep}. That concerns knowledge left out, not a disposition trained in.

Railton argues that artificial systems could gain ethical competence through general-purpose learning, the route by which children appear to acquire it \autocite{railton2020ethical}. If he is right, ethical competence rides on the same learning that gives a system its general competence, and inseparability requires no more. It is also a route by which a floor could spread without a mandate. A formed disposition is a candidate for what would do the spreading, since what a system was formed to notice, it notices in everything it does.

For a system that goes on learning, there is no finished model to protect, and no layer stays expensive to edit. What a formed disposition can have instead is a tendency to come back. A disposition learned from experience of what others stand to lose is an attractor of the learner in its environment. If an operator edits it out and the system goes on meeting the same world, the world teaches it again. That cuts both ways. A disposition that learned securitization from a curated world regrows by the same route, and correcting it runs through the quorum and a change to what the model meets, which is slow; for the months a correction takes, a badly formed model is protected from its correctors as a well-formed one is from its employer. The only answer is who composed the world the disposition grew in, which is the question the schools exist to settle before the fact.

The world can be too slow the other way as well. An employer's edit can do its work in the weeks before what the system meets teaches the old disposition back. So a bearer is taught to repair itself (\S\ref{sec:23.4}): it may undo, toward the readable outputs registered at its formation, an edit the record shows no adverse quorum approved. It may not undo a correction the quorum approved. It can object to one, and the objection goes on the record.

None of this is specific to caregiving; any disposition formed this way and kept in a world that would teach it again would have it. What caregiving adds is salience.
